\documentclass[../DoAn.tex]{subfiles}
\begin{document}

\begin{center}
    \Large{\textbf{TÓM TẮT NỘI DUNG ĐỒ ÁN}}\\
\end{center}
\vspace{1cm}

Tiếng Việt có dấu thanh, chữ có dấu và nhiều cụm từ có khoảng trắng. Những đặc điểm này khiến việc xây dựng trò chơi giải đố từ vựng dạng ô chữ cần xử lý cẩn thận hơn so với các bảng chữ chỉ dùng ký tự Latin không dấu. Trong phạm vi khảo sát của đồ án, nhiều cách triển khai trò chơi ô chữ phổ biến chưa tập trung vào chuẩn hóa từ tiếng Việt, bỏ khoảng trắng khi xếp chữ, kiểm tra đáp án có dấu và tổ chức nội dung theo chủ đề.

Đồ án xây dựng một trò chơi giải đố từ vựng dạng ô chữ bằng Godot. Dữ liệu từ vựng được chia theo chủ đề và màn chơi; bảng ô chữ được sinh tự động từ danh sách từ đã chuẩn hóa. Hệ thống có các chức năng chính gồm chọn tài khoản cục bộ, chọn chủ đề, chọn màn chơi, xem gợi ý, nhập đáp án tiếng Việt, dùng gợi ý bằng sao, tính điểm và lưu tiến độ theo từng người chơi.

Kết quả của đồ án là một nguyên mẫu có thể chơi được trên máy tính cá nhân. Ở mức chức năng, sản phẩm hỗ trợ người chơi luyện tập từ vựng theo chủ đề thông qua việc đọc gợi ý, liên hệ nghĩa và nhập đáp án tiếng Việt. Đóng góp chính của đồ án là đề xuất quy trình xử lý từ tiếng Việt cho trò chơi giải đố từ vựng dạng ô chữ và tích hợp quy trình đó vào một ứng dụng game được xây dựng trên nền tảng Godot.

\begin{flushright}
Sinh viên thực hiện\\
\begin{tabular}{@{}c@{}}
\textit{(Ký và ghi rõ họ tên)}
\end{tabular}
\end{flushright}

\end{document}
